% Figure: Gerschgorin discs in the complex plane for a 4x4 matrix,
% with the true eigenvalues overlaid. Two discs are isolated and each
% holds exactly one eigenvalue; two overlap and jointly hold two.
\begin{figure}[htbp]
\centering
\begin{tikzpicture}
\begin{axis}[
  width=12cm,
  height=8cm,
  axis equal image,
  xlabel={$\operatorname{Re}$},
  ylabel={$\operatorname{Im}$},
  xmin=-3.2, xmax=11.5,
  ymin=-3.6, ymax=3.6,
  axis lines=middle,
  tick label style={font=\footnotesize},
  label style={font=\small},
  title={Gerschgorin discs and the true spectrum},
  title style={font=\bfseries},
]
% Disc 1: centre 0, radius 1 (isolated).
\draw[draw=engblue, fill=engblue!12] (axis cs:0,0) circle [radius=1];
% Disc 2: centre 3, radius 1.5 (overlaps disc 3).
\draw[draw=engblue, fill=engblue!12] (axis cs:3,0) circle [radius=1.5];
% Disc 3: centre 4, radius 1 (overlaps disc 2).
\draw[draw=engblue, fill=engblue!12] (axis cs:4,0) circle [radius=1];
% Disc 4: centre 9, radius 2 (isolated).
\draw[draw=engblue, fill=engblue!12] (axis cs:9,0) circle [radius=2];
% Disc centres.
\addplot[only marks, mark=x, mark size=3pt, engblue, thick] coordinates {
  (0,0) (3,0) (4,0) (9,0)
};
% True eigenvalues (illustrative): one in disc 1, two in the merged
% disc 2+3, one in disc 4.
\addplot[only marks, mark=*, mark size=2.4pt, engred] coordinates {
  (-0.4,0.3) (2.6,0.8) (4.3,-0.6) (8.7,0.0)
};
\node[font=\footnotesize, engblue, anchor=south] at (axis cs:0,1.05) {$D_{1}$};
\node[font=\footnotesize, engblue, anchor=south] at (axis cs:3,1.55) {$D_{2}$};
\node[font=\footnotesize, engblue, anchor=south] at (axis cs:4,1.05) {$D_{3}$};
\node[font=\footnotesize, engblue, anchor=south] at (axis cs:9,2.05) {$D_{4}$};
\node[font=\footnotesize, engred, anchor=west] at (axis cs:9.3,0.0)
  {true $\lambda$};
\end{axis}
\end{tikzpicture}
\caption[Gerschgorin discs localise the spectrum from the
entries]{Four Gerschgorin discs (blue) for a $4 \times 4$ matrix,
each centred at a diagonal entry $A_{ii}$ with radius equal to the
absolute row sum of the off-diagonal entries. Every eigenvalue (red)
lies in at least one disc, the conclusion of the Gerschgorin circle
theorem, read off the entries without a single eigenvalue
computation. The first disc $D_{1}$ and the last $D_{4}$ are disjoint
from the rest, so each contains exactly one eigenvalue. The middle
discs $D_{2}$ and $D_{3}$ overlap; their union is disjoint from the
others and therefore contains exactly two eigenvalues, though the
theorem does not say which disc holds which. The picture is the
engineer's free first diagnostic: an isolated disc clear of the
origin certifies an eigenvalue of known sign and rough magnitude, and
a matrix all of whose discs avoid the origin is certified invertible
before any factorisation is run.}
\label{fig:vol02:ch10:gerschgorin-discs}
\end{figure}
